% !TeX program = pdflatex
\documentclass[12pt,a4paper]{article}

\newcommand{\notelanguage}{finnish}
% Shared preamble for course notes and exercise sets.

\usepackage[T1]{fontenc}
\usepackage{lmodern}
\usepackage[english]{babel}

\usepackage[a4paper,margin=30mm]{geometry}
\usepackage{microtype}
\usepackage{graphicx}
\usepackage{enumitem}

\usepackage{amsmath}
\usepackage{amssymb}
\usepackage{amsthm}
\usepackage{mathtools}
\usepackage{aliascnt}

\usepackage[hidelinks,pdfusetitle]{hyperref}

\numberwithin{equation}{section}
\setlist{itemsep=0.25em,topsep=0.5em}

\theoremstyle{plain}
\newtheorem{theorem}{Theorem}[section]
\newaliascnt{lemma}{theorem}
\newtheorem{lemma}[lemma]{Lemma}
\aliascntresetthe{lemma}
\newaliascnt{proposition}{theorem}
\newtheorem{proposition}[proposition]{Proposition}
\aliascntresetthe{proposition}
\newaliascnt{corollary}{theorem}
\newtheorem{corollary}[corollary]{Corollary}
\aliascntresetthe{corollary}

\theoremstyle{definition}
\newaliascnt{definition}{theorem}
\newtheorem{definition}[definition]{Definition}
\aliascntresetthe{definition}
\newaliascnt{example}{theorem}
\newtheorem{example}[example]{Example}
\aliascntresetthe{example}
\newaliascnt{problem}{theorem}
\newtheorem{problem}[problem]{Problem}
\aliascntresetthe{problem}

\theoremstyle{remark}
\newaliascnt{remark}{theorem}
\newtheorem{remark}[remark]{Remark}
\aliascntresetthe{remark}

\usepackage[nameinlink,noabbrev]{cleveref}

\crefname{theorem}{theorem}{theorems}
\crefname{lemma}{lemma}{lemmas}
\crefname{proposition}{proposition}{propositions}
\crefname{corollary}{corollary}{corollaries}
\crefname{definition}{definition}{definitions}
\crefname{example}{example}{examples}
\crefname{problem}{problem}{problems}
\crefname{remark}{remark}{remarks}

\DeclareMathOperator{\im}{im}
\DeclareMathOperator{\rank}{rank}
\DeclarePairedDelimiter{\abs}{\lvert}{\rvert}
\DeclarePairedDelimiter{\norm}{\lVert}{\rVert}
\DeclarePairedDelimiter{\set}{\lbrace}{\rbrace}


\title{Harjoituksen 1 ratkaisut}
\author{Joona Piirainen}
\date{Syksy 2026}

\begin{document}

\maketitle

\noindent
Tehtävänannot ovat
\href{https://moodle.helsinki.fi/pluginfile.php/6978847/mod_resource/content/5/harj1.pdf}{Moodle-järjestelmässä}.

\section*{Ratkaisut}

\begin{exercise}
  Osoita induktiolla, että
  \[
    1^3 + 2^3 + 3^3 + \dots + n^3 =\left(\frac{n(n + 1)}{2}\right)^2
  \]
  kaikilla \(n \in \mathbb{N}\), \(n \geq 1\).
\end{exercise}

\begin{solution}
  Osoitetaan lause todeksi induktiolla. Osoitetaan ensiksi, että lause
  pitää paikkaansa, kun \(n = 1\). Saamme
  \[
    1^3 = \left(\frac{1(1 + 1)}{2}\right)^2 = 1.
  \]
  Oletetaan, että lause on tosi jollain kokonaisluvulla \(n \geq 1\). Sitten
  \begin{align*}
    1^3 + 2^3 + \cdots + n^3 + (n + 1)^3
    &= (1^3 + 2^3 + \cdots + n^3) + (n + 1)^3 \\
    &\overset{\text{I.O.}}{=} \left(\frac{n(n + 1)}{2}\right)^2
    + (n + 1)^3 \\
    &= (n + 1)^2\left(\frac{n^2}{4}\right) + (n + 1)(n + 1)^2 \\
    &= (n + 1)^2\left(\frac{n^2}{4} + n + 1\right) \\
    &= \frac{(n + 1)^2(n^2 + 4n + 4)}{4} \\
    &= \frac{(n + 1)^2(n + 2)^2}{4} \\
    &= \left(\frac{(n + 1)((n + 1) + 1)}{2}\right)^2
  \end{align*}
  Siis lause pätee myös luvulla \(n + 1\), joten induktioperiaatteen nojalla
  se pätee kaikilla kokonaisluvuilla \(n \geq 1\).
\end{solution}

\begin{exercise}
  Osoita induktiolla, että jos \(n \in \mathbb{N}\), niin \(n^3 - n\)
  on jaollinen kuudella.
\end{exercise}
\begin{solution}
  Osoitetaan lause todeksi induktiolla. Osoitetaan ensin, että lause pitää
  paikkaansa, kun \(n = 0\). Saamme \(0^3 - 0 = 0 = 6\cdot0\), joten tapaus
  \(n = 0\) on osoitettu.

  Olkoon \(n \geq 0\) mielivaltainen ja oletetaan, että \(6 \mid (n^3 - n)\).
  Osoitetaan
  \[
    6 \mid \left((n + 1)^3 - (n + 1)\right),
  \]
  eli, että \((n + 1)^3 - (n + 1)\) on kuudella jaollinen. Sieventämällä saadaan
  \begin{align}
    (n + 1)^3 - (n + 1) &= n^3 + 3n^2 + 3n + 1 - (n + 1) \notag \\
    &= n^3 + 3n^2 + 2n \notag \\
    &= (n^3 - n) + (3n^2 + 3n) \notag \\
    &= (n^3 - n) + 3n(n + 1). \label{eq:exercise-2-induction-step}
  \end{align}
  Induktio-oletuksen perusteella \(6 \mid (n^3 - n)\). Lisäksi joko \(n\) tai
  \(n + 1\) on parillinen, joten \(2 \mid n(n + 1)\). On siis olemassa
  \(k,m \in \mathbb{Z}\), joille on voimassa
  \(6k \overset{\text{I.O.}}{=} n^3 - n\), \(2m = n(n + 1)\).
  Näistä jälkimmäisestä seuraa \(3n(n + 1) = 3(2m) = 6m\), joten \(6
  \mid 3n(n + 1)\).
  Näin ollen yhtälön \eqref{eq:exercise-2-induction-step} oikean puolen
  molemmat termit ovat kuudella jaollisia, koska
  laskemalla yhtälöt \(6k = n^3 - n\) ja \(6m = 3n(n + 1)\) yhteen
  saadaan
  \[
    n^3 - n + 3n(n + 1) = 6k + 6m = 6(k + m),
  \]
  joten
  \begin{align}
    6 & \mid n^3 - n + 3n(n + 1) \\
    \implies
    6 & \mid \left((n + 1)^3 - (n + 1)\right).
  \end{align}
  Siis lause pätee myös luvulla \(n + 1\), joten induktioperiaatteen
  nojalla se pätee kaikilla luonnollisilla luvuilla.
\end{solution}

\begin{exercise}
  Osoita induktiolla, että
  \[
    1 \cdot 1! + 2 \cdot 2! + 3 \cdot 3! + \dots + n \cdot n! = (n + 1)! - 1
  \]
  kaikilla \(n \in \mathbb{N}, n \geq 1\).
\end{exercise}

\begin{solution}
  Osoitetaan ensin, että lause on tosi, kun \(n = 1\). On voimassa
  \(1 \cdot 1! = 1 = 2! - 1 = (1 + 1)! - 1\).

  Olkoon \(n \geq 1\) mielivaltainen ja oletetaan, että
  \[
    1 \cdot 1! + 2 \cdot 2! + \cdots + n \cdot n! = (n + 1)! - 1.
  \]
  Osoitetaan
  \[
    1 \cdot 1! + 2 \cdot 2! + \cdots + n \cdot n!
    + (n + 1)(n + 1)! = (n + 2)! - 1.
  \]
  Sieventämällä saadaan
  \begin{align*}
    &\phantom{{}={}} 1 \cdot 1! + 2 \cdot 2! + \cdots + n \cdot n!
    + (n + 1)(n + 1)! \\
    &= \left(1 \cdot 1! + 2 \cdot 2! + \cdots + n \cdot n!\right)
    + (n + 1)(n + 1)! \\
    &\overset{\text{I.O.}}{=} (n + 1)! - 1 + (n + 1)(n + 1)! \\
    &= 1 \cdot (n + 1)! + (n + 1)!(n + 1) - 1 \\
    &= \left(1 + (n + 1)\right)(n + 1)! - 1 \\
    &= (n + 2)(n + 1)! - 1 \\
    &= (n + 2)! - 1.
  \end{align*}
  Siis lause pätee myös luvulla \(n + 1\), joten induktioperiaatteen
  nojalla se pätee kaikilla \(n \in \mathbb{N}, n \geq 1\).
\end{solution}

\begin{exercise}
  Me ``todistamme'', että kaikki koirat ovat samanvärisiä induktion avulla.
  Tarkastellaan ensin yhtä koiraa. Se on selvästi samanvärinen itsensä kanssa,
  joten ryhmässä, jossa on yksi koira, pätee että kaikki koirat ovat
  samanvärisiä. Teemme induktio-oletuksen, että \(n = k\) koiraa on
  samanvärisiä. Todistetaan, että silloin myös \(k + 1\) koiraa on
  samanvärisiä. Tarkastellaan ryhmää, jossa on \(k + 1\) koiraa. Tästä
  ryhmästä voidaan poistaa yksi koira, jolloin jäljelle jää \(k\) koiraa.
  Induktio-oletuksen mukaan nämä \(k\) koiraa ovat samanvärisiä. Palautetaan
  poistettu koira alkuperäiseen ryhmään ja poistetaan siitä jokin toinen
  koira. Jälleen jäljelle jää \(k\) koiraa ja niiden täytyy olla
  samanvärisiä. Mutta nyt myös ryhmästä poistetut koirat täytyy olla
  samanvärisiä näiden kanssa, ja siten kaikkien \(k + 1\) koiran täytyy olla
  samanvärisiä. Tästä seuraa, että kaikkien koirien täytyy olla samanvärisiä.
  Missä on virhe?
\end{exercise}

\begin{solution}
  Merkitään väitettä symboleilla \(P(k)\): jokaisessa \(k\) koiran ryhmässä
  kaikki koirat ovat samanvärisiä. Perustapaus \(P(1)\) on tosi. Väitetty
  induktioaskel ei kuitenkaan päde tapauksessa \(k = 1\). Kun kahden koiran
  ryhmästä poistetaan vuorotellen toinen koira, saadaan kaksi yhden koiran
  ryhmää, joiden leikkaus on tyhjä. Ryhmissä ei siis ole yhteistä koiraa,
  jonka avulla koirien voitaisiin päätellä olevan samanvärisiä. Näin ollen
  johtopäätöstä \(P(1) \implies P(2)\) ei voida tehdä.

  Todistuksen oikeellisuus vaatisi induktioaskeleen pitävyyden kaikilla
  \(k \in \mathbb{N}, k \geq 1\), eli tulisi olla voimassa
  \begin{equation}\label{eq:exercise-4-induction-step}
    P(k) \implies P(k + 1)
    \qquad \text{kaikilla } k \in \mathbb{N},\ k \geq 1.
  \end{equation}
  Edellä on kuitenkin perusteltu, että \eqref{eq:exercise-4-induction-step} ei
  ole voimassa kun \(k = 1\).

  Todetaan vielä, että induktioaskel pätee, kun \(k \geq 2\).
  Poistamalla \(k + 1\) koiran ryhmästä vuorotellen kaksi eri koiraa saadaan
  kaksi \(k\) koiran ryhmää. Induktio-oletuksen mukaan kumpikin ryhmä on
  samanvärinen. Ryhmien leikkauksessa on \(k - 1 \geq 1\) koiraa, joten niillä
  on yhteinen koira ja siten sama väri. Ainoa puuttuva induktioaskel on siis
  \(P(1) \implies P(2)\).
\end{solution}
\end{document}
